\chapter {Introduction}
\label{label}


The introduction is the first section of the document and covers the 30\% of the whole dissertation. It draws your reader in, setting the stage for your research with a clear focus, purpose, and direction on a relevant topic.
It may include a review of the background theory in your subject area in the form of a literature review. The literature review may be in a separate section. ~\cite{BernersLee2010}. ~\cite{Abdullah2017, Abidi2017, Angelov2019, Auxier2021, Bauman2004, Bramer2007, Ek2018}

The introduction should include:

•	Dissertation topic, in context: give necessary background information

•	Your focus and scope: focus and define the scope of your research

•	The relevance of your research: define how your research relates to existing work on your topic. Explain how it solves a practical or theoretical problem, addresses a gap in the literature, builds on existing research and proposes a new understanding of the topic. 

•	Your questions and objectives: explain what your research aims to find out, and how.

•	An overview of your structure: end the introduction with an outline of the structure of your dissertation to follow. Share a brief summary of each chapter and show clearly how each contributes to your central aims.  


\section{Aims and Objectives}
\label{label1}


\section{Research Questions}
\label{label2}


\section{Contribution}
\label{label3}


\section{Structure of the Thesis}
\label{label4}


\section{Summary}
\label{label5}